\clearpage\thispagestyle{empty}
\pdfbookmark[0]{Quellen zum Kometen (18. September 1815)}{comet1815sep-sources}
\begin{center}\textbf{Quellen zum Kometen (18. September 1815)}\end{center}
\noindent Diese KI-Übertragung ins heutige Deutsch folgt dem Bericht in den \href{https://gdz.sub.uni-goettingen.de/id/PPN236007467}{\textit{Werken}, Band~VI} (Göttingen, 1874), S.~389--391. Der Erstabdruck erschien in den \textit{Göttingischen gelehrten Anzeigen}, Stück~149 (18. September 1815), S.~\href{https://api.digitale-sammlungen.de/iiif/image/v2/bsb10538688_00113/full/1600,/0/default.jpg}{1473}--\href{https://api.digitale-sammlungen.de/iiif/image/v2/bsb10538688_00116/full/1600,/0/default.jpg}{1476}. Die Links zum Erstabdruck führen zum Exemplar der Bayerischen Staatsbibliothek. Der Verweis auf Möbius’ Elemente führt zu \href{https://api.digitale-sammlungen.de/iiif/image/v2/bsb10538617_00593/full/1600,/0/default.jpg}{\textit{Monatliche Correspondenz}, Band~XXVIII, S.~577}, wo eine mit N. unterzeichnete Anmerkung Möbius’ Vergleich und seine verbesserten Elemente angibt.
\par\smallskip\noindent Im Erstabdruck steht auf S.~\href{https://api.digitale-sammlungen.de/iiif/image/v2/bsb10538688_00114/full/1600,/0/default.jpg}{1474} \textit{Fraunhofer}, und auf S.~\href{https://api.digitale-sammlungen.de/iiif/image/v2/bsb10538688_00115/full/1600,/0/default.jpg}{1475} stehen vor dem Abschnitt über Juno drei Sternchen. Beide Drucke geben für die Rektaszension des Kometen am 30.~März 59~13~3 und für die Deklination der Juno am 29.~März 30.1 Sekunden an (im Erstabdruck 30,1). Im gesammelten Augustbrief an Bode stehen dagegen 59~13~2 und 38.1. Beide Drucke des Berichts haben beim Zeiteintrag für den 1.~März ein Doppelstrichzeichen.
\par\smallskip\noindent Im Erstabdruck steht bei der Perihellänge der Juno 53~14~53,8 ohne Einheitenzeichen und bei der Bahnneigung $13^\circ\,14'\,0''1$. In der Werkausgabe ist die Länge mit Einheitenzeichen versehen, die Neigung dagegen ohne solche Zeichen gesetzt. Der letzte Satz nennt im Erstabdruck \textit{Hrn.\ N.}, in der Werkausgabe Nicolai. Die Übertragung folgt den Lesarten der Werkausgabe. KI-Übertragung und KI-Satz.
\clearpage
